% !PRINTED 385
\weberpage{385}

\weq{5}{\omega_n' + 1 = \frac{1}{Q_n}\left\{ Q_n - Q_{n-1} - \frac{(-1)^n}{Q_n\left(\omega' - \dfrac{P_n}{Q_n}\right)} \right\}.}

If now $n$ is taken sufficiently large, then
\[
\omega' - \frac{P_n}{Q_n}, \quad \omega' - \frac{P_{n-1}}{Q_{n-1}}
\]
% [sic] The source prints $\omega'-\omega=2y\sqrt d$; from (1)--(2) of §122
% the left side is actually $-2y\sqrt d$.  The proof here uses only non-vanishing.
differ arbitrarily little from the non-vanishing difference $\omega' - \omega = 2y\sqrt{d}$, whence (4) shows that $\omega_n'$ is negative. But from (5) it follows that $\omega_n' + 1$ is positive, since $Q_n - Q_{n-1}$, being a positive integer, is at least equal to 1, and
\[
\frac{1}{Q_n\left(\omega' - \dfrac{P_n}{Q_n}\right)}
\]
becomes arbitrarily small for a sufficiently large $n$; and with this the proof of Theorem 1. is complete.

Since the conjugate number of the conjugate $\omega'$ is the original number $\omega$, it follows from the definition of the reduced numbers \gesperrt{that $\omega$ and $-1:\omega'$ are reduced simultaneously.}

\gesperrt{2. If $\omega$ is reduced and, for an integer $\alpha$,}

\weq{6}{\omega = \alpha + \frac{1}{\omega_1},}

\gesperrt{then $\omega_1$ is reduced if and only if $\alpha$ is the greatest integer contained in $\omega$, that is, if}

\weq{7}{\alpha < \omega < \alpha + 1.}

If the condition (7) is fulfilled, then $\omega_1$ is positive and greater than 1, and

\weq{8}{\omega_1' = \frac{1}{\omega' - \alpha}}

is, since $\omega'$ is negative, a negative proper fraction; hence $\omega_1$ is reduced. On the other hand, $\omega_1$ cannot be reduced if $\alpha > \omega$, since then $\omega_1$ would be negative, or if $\alpha + 1 < \omega$, since otherwise $\omega_1$ would be less than 1.

If one puts (8) in the form

\weq{9}{-\frac{1}{\omega_1'} = \alpha + \frac{1}{-\dfrac{1}{\omega'}},}
%CONT>>
